\documentclass[12pt]{article}
\input{bayesuvius.sty}
\usepackage{authblk}
\title{Determination of Drought Triggering Probability \\via Bayesian Network}
\author[1]{Nishadh Kalladath}
\author[1,2]{Robert R. Tucci}
\affil[1]{ICPAC (IGAD Climate Prediction and Applications Centre)}
\affil[2]{www.ar-tiste.xyz}
\begin{document}
\maketitle
\abstract{xxx}
\section{Introduction}

{\bf Drought Anticipatory Action} (DAA)

bnet (Bayesian network) (see Ref.\cite{Bayesuvius})

{\bf Standardized Precipitation Index} (SPI) ranges from $-4$ (least rain) to $4$ (most rain).

\begin{figure}[h!]
\centering
\includegraphics[width=6in]
{tim-palmer-cartoon.png}
\caption{Illustration
from Tim Palmer's book Ref.\cite{palmer-book}}
\label{fig-tim-palmer-cartoon}
\end{figure}


\section{Notation
and Conventions}
We will use $\indi(S)$ 
to denote the indicator
function
(a.k.a. the truth function).
$\indi(S)$ equals 1 if statement $S$ is true,
and 0 if $S$ is false.
For example, $\indi(x\geq 0)$
equals 0 if $x<0$ and 1 otherwise.

In this article, random variables will
be indicated by underlining.
For example, $P(\rvx>x)$
will denote
the probability 
that the random 
variable $\rvx$
is greater than $x$.


Let

$N=$ the number of times (seasons) 
in which predicted and observed droughts were compared.

$N_{x|a}=$  number of events in which a number of 
$a\in \bool$ droughts occurred and 
a number of $x\in \bool$ droughts were predicted.
For example, $N_{1|0}$ is the number of events 
in which no drought actually occurred but 
it was predicted nonetheless (False Alarm).


$N_{0|0}=$ {\bf True Negatives} (TN), correct rejection

$N_{1|1}=$ {\bf True Positives} (TP), hit

$N_{0|1}=$ {\bf False Negatives} (FN), miss, type II error or underestimation

$N_{1|0}=$ {\bf False Positives} (FP), false alarm, type I error or overestimation



\beq
N_{|a} = \sum_{x'} N_{x'|a},
\quad
N_{x|} = \sum_{a'} N_{x|a'}
\eeq

\beq
N = \sum_x \sum_a N_{x|a} =
\sum_a N_{|a}
=\sum_x N_{x|a}
\eeq

\beq
P(x|a) = R_{x|a}= \frac{N_{x|a}}{N_{|a}}
,\;
P'(a|x) = R'_{a|x}= \frac{N_{x|a}}{N_{x|}}
\eeq

\beq
\sum_x P(x|a) =
\sum_a P'(a|x)=1
\eeq
Note that both 
$P(x|a)$
and $P'(a|x)$ 
are
conditional
probability distributions.

Note that $N_{x|a}$
for $x,a\in \bool$
yields 4 degrees of
freedom (dofs).
Instead of those 4  dofs,
we will use $N$
and 3 others dofs,
namely,
hit rate (HR), 
false alarm rate (FAR)
and area under ROC curve (AUROC).
HR, FAR and AUROC
are defined
in Appendix \ref{app-n-ratios}




\section{DAA bnet}


Let

$\kappa=(\sigma, \delta)=$
the drought kind, 
where  $\sigma\in\{\text{MAM, JJAS}\}$\footnote{Letters here stand for the first letter of the months of the year.} is the season, and $\delta\in\{\text{Moderate, Extreme, Severe}\}$
is the drought degree.

$I_\kappa=$ {\bf Impact } (i.e., {\bf severity}) of drought is defined as minus SPI (higher rain,
higher SPI, less 
severe drought)

$ P(I_\kappa)=$ likelihood, 
probability  that a drought of severity $I_\kappa$  will occur in the near future (weeks to months), based on forecasts.

$\mu=$ ensemble member index

$M=$ number of ensemble members, 51 for $ECMWF$ (European Centre for Medium-Range Weather Forecasts)


\beqa
Risk_\kappa &=& E[\rvI_\kappa] = \frac{1}{M} \sum_{\mu=1}^{M} I_{\kappa,\mu}
\\
&=&\sum_{I_\kappa} P(I_\kappa) I_\kappa
\eeqa

A {\bf risk matrix} is
a  $4\times 4$ matrix
with $X$ axis 
equal to impact ($I$)
and $Y$ axis equal to likelihood
$P$, with 
$I$ binned into four values
and $P$ binned into 4 values.
The matrix elements give
the value of
some arbitrary function 
$f(I, P)$ defined as risk. Usually, $f(I, P)=IP$ but not necessarily. I find 
risk matrices arbitrary
and not very useful. The 
Wikipedia article on risk matrices (Ref.\cite{wiki-risk-matrix}) agrees.

\begin{table}[!h]
\begin{tabular}{|l|l|l|l|}
\hline
$\sigma\backslash\delta$ & \cellcolor[HTML]{FFFFC7}Moderate & \cellcolor[HTML]{FFFFC7}Extreme & \cellcolor[HTML]{FFFFC7}Severe \\ \hline
\cellcolor[HTML]{FFFFC7}MAM & $-$0.55 & $-$0.98 & $-$0.99 \\ \hline
\cellcolor[HTML]{FFFFC7}JJAS & $-$0.41 & $-$0.98 & $-$0.99 \\ \hline
\end{tabular}
\caption{SPI threshold  for Karamoja, Uganda.
Cells give $-I^*_\kappa$,
where $\kappa=(\sigma, \delta)$}
\label{tab-spi-karamoja}
\end{table}

$I_\kappa^*  =$ impact threshold 
   for drought (See Table
   \ref{tab-spi-karamoja}
   for examples
   of $I^*_\kappa$.)
   

\beq
P(\rvI_\kappa\geq I_\kappa^*  ) = \frac{1}{M} \sum_{\mu=1}^{M} \indi(I_{\kappa,\mu} \geq I_\kappa^*) 
\eeq

Let $X^i_\kappa\in \bool$ for $i=1,2, 3, \ldots, nx$ represent vulnerability to local problems
such as high road density, high population density,
historical 
catastrophic events, etc.
It equals 1 if a particular danger is present, and 0 if not.

Let

%\beq
%N_\kappa = \sum_{i=1}^{nx}
%X^i_\kappa
%\eeq

$\eps \in \RR$ be a normally distributed random variable
with mean zero and
some a priori
given deviation.

$RC_\kappa^i=$ replacement cost for 
vulnerability $X^i_\kappa$

$IC_\kappa^i=$  insurance (i.e., preventive) cost for vulnerability $X^i_\kappa$


Note that insurance always costs
more than replacement, so

\beq
0\leq IC_\kappa^i \leq RC_\kappa^i
\eeq

$P_{danger, \kappa}^i=
\frac{IC_\kappa^i}{RC_\kappa^i}$ probability of danger for Vulnerability $X_\kappa^i$.

Note that less 
danger 
means less 
insurance cost.


$P^*_\kappa=$ probability threshold for drought

\beq
P^*_\kappa =
f_\kappa(X^1_\kappa, X^2_\kappa,
\ldots, X^{nx}_\kappa) + \eps 
\eeq
where we model the function $f_\kappa()$ by

\beq
f_\kappa(X^1_\kappa, X^2_\kappa,
\ldots, X^{nx}_\kappa) 
=
1 -
\max_i\{P^i_{danger, \kappa}
X_\kappa^i\}
\eeq



$TD_\kappa=$ triggering decision to do drought warning

$P^*_\kappa=$ threshold probability that triggers a  drought warning



\beq
TD_\kappa=
\indi\left[P(\rvI_\kappa\geq I_\kappa^*  ) > P^*_\kappa\right]
\eeq



\begin{figure}[h!]
$$
\xymatrix{
\rvX^1_\kappa \ar[dr]
&\rvX^2_\kappa \ar[d]
&\rvX^3_\kappa\ar[dl]
\\
\rvP(\rvI_\kappa> I^*_\kappa)\ar[d] 
& \rvP^*_\kappa\ar[dl]
&\rv{\eps}\ar[l]
\\
\rv{TD}_\kappa
\ar[d]\ar[dr]\ar[drr]
\\
\ul{FAR}_\kappa 
&\ul{HR}_\kappa
&\ul{AUROC}_\kappa
}
$$
\caption{DAA bnet.
We use one of these 
bnets for each 
$\kappa$ value and for each grid cell
of 0.25 degrees in latitude by 0.25 degrees 
in longitude. 1 degree in latitude
or longitude is approximately 100 km.
}
\label{fig-aa-bnet}
\end{figure}


The TPMs\footnote{TPM=Transition
Probability Matrix.
This is the same thing
as CPT=Conditional Probability Table
} printed in blue, for the
bnet Fig.\ref{fig-aa-bnet}, are
as follows.\footnote{
A deterministic node satisfying
the structural equation $y=f(x_1, x_2)$,
has the TPM: $P(y|x_1, x_2)=
\delta(y, f(x_1, x_2))
$, where 
$\delta(a,b)=\indi(a=b)$ is the Kronecker delta function. }



\beq\color{blue}
P^*_\kappa =
f_\kappa(X^1_\kappa, X^2_\kappa,
\ldots, X^{nx}_\kappa) + \eps 
\eeq

\beq\color{blue}
TD_\kappa
=
\indi[
P(\rvI_\kappa>I^*_\kappa) >P^*_\kappa]
\eeq

\beq\color{blue}
P(Y_\kappa^i|TD_\kappa)
= \text{to be learned empirically from dataset}
\eeq
where $(Y_\kappa^1, Y_\kappa^2, Y_\kappa^3) =
(FAR_\kappa, HR_\kappa,
AUROC_\kappa)$.

\begin{mdframed}[hidealllines=true,backgroundcolor=blue!10]
\noindent {\bf Algorithm for finding best triggers $T^*_\kappa$:}

Enter bnet Fig.\ref{fig-aa-bnet} into bnet software
PyAgrum
(Ref.\cite{pyagrum}).
Enter into software available
evidence for
nodes $P(\rvI_\kappa> I^*_\kappa)$, $(X^1, X^2, X^3)$
and $(Y_\kappa^1, Y_\kappa^2,Y_\kappa^3)$.
Allow software
to
infer the 
distribution of 
$P^*_\kappa$.
Choose for
$P^*_\kappa$
the state
with maximum 
probability. 
\end{mdframed}



\appendix 

\section{Appendix: Functions of $N_{x|a}$}
\label{app-n-ratios}

Terminology Table
Adapted from Wikipedia Ref.\cite{wiki-roc}

positive events (P):
$P=N_{|1}=\sum_x N_{x|1}$,
the number of real positive 
events in the data

negative events (N):
$N=N_{|0}=\sum_x N_{x|0}$,
the number of real negative events in the data


sensitivity, recall, true positive rate (TPR),
or hit rate  ({\color{red}HR}):
\beq TPR=
{  {R_{1|1}} ={\frac { {N_{1|1}} }{ {N_{|1}} }}=
{\frac { {N_{1|1}} }{ {N_{1|1}} + {N_{0|1}} }}=1- {R_{0|1}} }
\eeq

specificity, selectivity or true negative rate (TNR):
\beq TNR=
{  {R_{0|0}} ={\frac { {N_{0|0}} }{ {N_{|0}} }}
={\frac { {N_{0|0}} }{ {N_{0|0}} + {N_{1|0}} }}=
1- {R_{1|0}} }
\eeq

precision or positive predictive value (PPV):
\beq 
{  {PPV} ={\frac { {N_{1|1}} }
{ {N_{1|1}} + {N_{1|0}} }}=1- {FDR} }
\eeq

negative predictive value (NPV):
\beq 
{  {NPV} ={\frac { {N_{0|0}} }
{ {N_{0|0}} + {N_{0|1}} }}=1- {FOR} }
\eeq

miss rate or false negative rate (FNR):
\beq FNR=
{  {R_{0|1}} ={\frac { {N_{0|1}} }
{ {N_{|1}} }}={\frac { {N_{0|1}} }{ {N_{0|1}} 
+ {N_{1|1}} }}=1- {R_{1|1}} }
\eeq

fall-out, false positive rate (FPR),
or false alarm rate ({\color{red}FAR}):
\beq FPR=
{  {R_{1|0}} ={\frac { {N_{1|0}} }
{ {N_{|0}} }}={\frac { {N_{1|0}} }
{ {N_{1|0}} + {N_{0|0}} }}=1- {R_{0|0}} }
\eeq

false discovery rate (FDR):
\beq 
{  {FDR} ={\frac { {N_{1|0}} }
{ {N_{1|0}} + {N_{1|1}} }}=1- {PPV} }
\eeq

false omission rate (FOR):
\beq 
{  {FOR} ={\frac { {N_{0|1}} }
{ {N_{0|1}} + {N_{0|0}} }}=1- {NPV} }
\eeq


accuracy (ACC):
\beq 
{  {ACC} ={\frac { {N_{1|1}} + {N_{0|0}} }
{ {N_{|1}} + {N_{|0}} }}={\frac { {N_{1|1}} + {N_{0|0}} }
{ {N_{1|1}} + {N_{0|0}} + {N_{1|0}} + {N_{0|1}}
 }}}
\eeq

balanced accuracy (BA):
\beq 
{  {BA} ={\frac {R_{1|1}+R_{0|0}}{2}}}
\eeq

F1 score
is the harmonic mean of precision and sensitivity: 
\beq 
{  {F}_{1}=2\times {\frac { {PPV} 
\times  {R_{1|1}} }{ {PPV} + {R_{1|1}} }}=
{\frac {2 {N_{1|1}} }{2 {N_{1|1}} + {N_{1|0}} +
 {N_{0|1}} }}}
\eeq

The area under the ROC (Receiver
Operating Characteristic) curve,
known as AUC or {\color{red}AUROC},
is defined by 
Fig.\ref{fig-roc-auc}. One set of 4
values 
$N_{x|a}$
where $a,x\in \bool$,
only gives one
point of the ROC curve. To
get a full ROC curve,
we need to vary a parameter. In the
case of DAA,
we vary geographical location,
for example.

\begin{figure}[h!]
\centering
\includegraphics[width=2in]
{roc-auc.png}
\caption{Example of ROC.
Green shaded
 area is the AUC of the ROC.} 
\label{fig-roc-auc}
\end{figure}

\bibliographystyle{plain}
\bibliography{aa_references}
\end{document}